\section*{第 \ref{ch:g10:soil-formation} 章——土壤：一种形成缓慢的资源}

\begin{solution}{exo:g10:soil-formation:1}
\omterm{def:g10:soil-formation:chemical-weathering}{化学风化}是指岩石中的矿物通过
与水及其携带的气体（二氧化碳、氧气）反应而转变：
一些矿物溶解，另一些变为在地表稳定的新矿物。
\omterm{def:g10:soil-formation:hydrolysis}{水解}是指硅酸盐与水及其氢离子
发生的\omterm{def:g10:soil-formation:chemical-weathering}{化学风化}：金属离子进入溶液，矿物的其余部分
重新构成为\omterm{def:g10:soil-formation:clay}{黏土矿物}。
\end{solution}

\begin{solution}{exo:g10:soil-formation:2}
\omterm{def:g10:soil-formation:weathering-profile}{腐岩}是原地分解的岩石：质软呈土状，保留了\omterm{def:g10:soil-formation:weathering-profile}{母岩}的颗粒、
条带和结构，并无间断地向下过渡到新鲜岩石中。
河沙上形成的土壤则覆盖在他处搬来的物质之上：
颗粒经过磨圆和分选，
完全没有其下方岩石的结构。
\end{solution}

\begin{solution}{exo:g10:soil-formation:3}
长石转变为黏土（高岭石），释放出钾、钠或
钙及二氧化硅进入水中；黑云母转变为黏土和
铁的氧化物；石英保持不变，以沙粒的形式留下。
\end{solution}

\begin{solution}{exo:g10:soil-formation:4}
暗色矿物中的铁，即 \ce{Fe^2+}，被水携带的氧气
氧化为不溶性的三价铁氧化物：黄褐色的针铁矿和红色的赤铁矿，
包裹颗粒表面，使整个物质呈现颜色。
\end{solution}

\begin{solution}{exo:g10:soil-formation:5}
两者分别为 \qty{0.017}{mm} 每年与 \qty{1.537}{mm} 每年，
后者约为前者的 $1.537/0.017 \approx 90$ 倍。
\end{solution}

\begin{solution}{exo:g10:soil-formation:6}
钙：1 与 1。碳：$1 + 1 = 2$ 与 2。氧：$3 + 2 + 1 = 6$ 与
$2 \times 3 = 6$。氢：2 与 2。电荷：0，且 $+2 - 2 = 0$。
\end{solution}

\begin{solution}{exo:g10:soil-formation:7}
方解石的溶解只产生溶解的离子 \ce{Ca^2+} 和
\ce{HCO3-}：全部被带走，纯石灰岩不留下固体残留物。
长石的\omterm{def:g10:soil-formation:hydrolysis}{水解}产生固体高岭石，它留存下来；
花岗岩中的石英几乎不发生反应。因此花岗岩
只留下黏土和沙。
\end{solution}

\begin{solution}{exo:g10:soil-formation:8}
黏土薄片和腐殖质带有负电荷，能吸住其表面带正电的
钾离子 \ce{K+}，使流经的雨水不会将其带走。根释放出
氢离子，它们取代薄片上的钾；
游离于土壤水中的钾
被根系吸收。
\end{solution}

\begin{solution}{exo:g10:soil-formation:9}
\omterm{def:g10:soil-formation:chemical-weathering}{化学风化}需要水，潮湿气候下水量充足，干燥气候下水量稀少
（寒冷气候下则冻结）；反应在温暖时进行得更快。
只有在高温和水量充足的条件下，经过长时间作用，
反应才能进行得足够彻底，几乎除尽所有可溶物质，
只留下铁和铝的氧化物。
\end{solution}

\begin{solution}{exo:g10:soil-formation:10}
形成 \qty{1}{cm} 厚的风化层需 $10/0.017 \approx 590$ 年；形成 \qty{30}{cm} 厚的风化层
需 $300/0.017 \approx 17\,600$ 年；形成 \qty{1}{m} 厚的风化层需 $1000/0.017 \approx 59\,000$ 年。
\end{solution}

\begin{solution}{exo:g10:soil-formation:11}
体积：$10\,000 \times 0.0015 = \qty{15}{m^3}$；质量：$15 \times 1200 =
\qty{18000}{kg}$，合每公顷每年 \num{18}~吨。
\end{solution}

\begin{solution}{exo:g10:soil-formation:12}
$46\,000/0.017 \approx 2.7$ 百万年。这样厚的\omterm{def:g10:soil-formation:weathering-profile}{腐岩}需要
地貌长期保持稳定，风化覆盖层未被剥蚀，
历时数百万年。
\end{solution}

\begin{solution}{exo:g10:soil-formation:13}
针铁矿：\ce{4 Fe^2+ + O2 + 6 H2O -> 4 FeO(OH) + 8 H+}：铁 4 与 4，
氧 $2 + 6 = 8$ 与 $4 \times 2 = 8$，氢 12 与 $4 + 8 = 12$，电荷
$+8$ 与 $+8$。赤铁矿：\ce{4 Fe^2+ + O2 + 4 H2O -> 2 Fe2O3 + 8 H+}：铁 4
与 4，氧 $2 + 4 = 6$ 与 6，氢 8 与 8，电荷 $+8$ 与 $+8$。
\end{solution}

\begin{solution}{exo:g10:soil-formation:14}
净损失为 $0.082 - 0.017 = \qty{0.065}{mm}$ 每年；$300/0.065 \approx 4600$
年。若采用常规耕作，$300/1.52 \approx 200$ 年，
所需时间缩短了二十多倍。
\end{solution}

\begin{solution}{exo:g10:soil-formation:15}
这一吨中含 \qty{300}{kg} 正长石。每千克可产生
\qty{464}{g} 高岭石：$300 \times 0.464 \approx \qty{139}{kg}$。钾：
$300 \times 0.1405 \approx \qty{42}{kg}$。
\end{solution}

\begin{solution}{pb:g10:soil-formation:1}
\textbf{1.} $4 \times 0.3048 \approx \qty{1.22}{m}$。

\textbf{2.} $0.5 \times 30.48 \approx \qty{15}{cm}$。

\textbf{3.} 腐殖质，即分解者由死去的
植物所形成的深色有机物。

\textbf{4.} 氧化铁，即针铁矿和赤铁矿，由
母质矿物中的二价铁经氧化形成。

\textbf{5.} 否。向下，剖面将过渡为风化岩，
即\omterm{def:g10:soil-formation:weathering-profile}{腐岩}（或其他风化物质），然后进入新鲜的
\omterm{def:g10:soil-formation:weathering-profile}{母岩}。

\textbf{6.} \ce{2 KAlSi3O8 + 2 H+ + 9 H2O -> Al2Si2O5(OH)4 + 2 K+ + 4 H4SiO4}。

\textbf{7.} $1000/278.33 \approx \qty{3.59}{mol}$。

\textbf{8.} $3.59/2 \approx \qty{1.80}{mol}$ 高岭石，其质量为 $1.80 \times 258.16
\approx \qty{464}{g}$。

\textbf{9.} 钾离子与溶解的二氧化硅：它们被
水带入河流，再由河流带入海洋。

\textbf{10.} 几乎不发生反应：石英颗粒被留下，成为沙。

\textbf{11.} $1219/0.017 \approx 72\,000$ 年。

\textbf{12.} $1.537 - 0.017 = \qty{1.52}{mm}$ 每年。

\textbf{13.} $1219/1.52 \approx 800$ 年。

\textbf{14.} $72\,000/800 = 90$：剖面流失的速度
是其形成速度的九十倍。

\textbf{15.} 否：$0.013 < 0.017$，侵蚀速度略慢于
生成速率，土壤能够维持自身。

\textbf{16.} $T = S/(E - P)$。

\textbf{17.} $300/1.52 \approx 200$ 年。

\textbf{18.} $300/(0.082 - 0.017) = 300/0.065 \approx 4600$ 年。

\textbf{19.} 土壤在几个世纪内流失，
需要数万年才能重新形成：以人的一生和社会的尺度而言，失去的
便无法得到替代。

\textbf{20.} $1000/1.52 \approx 660$ 年：\textbf{一米土壤将在约 660 年，
即几个世纪，在常规翻耕下流失。}
\end{solution}
